% ===== §15 Conclusions =====
\section{Conclusions: At the Boundaries of the Frameworks}
\label{sec:conclusions}

\noindent
This paper has traversed considerable ground, from the political economy of wealth theory to the semiotics of the happiness jar, from the psychoanalysis of phallic signifiers to the Hegelian-Marxist dialectic of relational wealth, from the feminist recovery of care labour to the epistemic justice of the right to know one's relational field. The traversal has been guided by a single question---what is wealth in intimate relations?---and has arrived at an answer that is at once theoretically precise and practically consequential: the primary form of wealth in intimate relations is Generative Relational Wealth, an emergent property of co-evolutionary dynamics that is stored in the geometry of the relational attractor landscape, inexhaustible under use, non-possessable by either party, and self-generating in its structure.

In keeping with the series' commitment to honest multi-framework analysis---the polyphonic conclusion form that refuses the false unity of a single synthetic statement---this section presents the conclusions that each of the paper's major frameworks can and cannot reach, ends with the three concentric circles that constitute the paper's structural closure, and states explicitly why no synthetic conclusion is offered.

\subsection{What the Political Economy Framework Can and Cannot Say}
\label{subsec:conc_pe}

\noindent
The political economy framework can say: that the major traditions of wealth theory share an individualist presupposition that renders GRW invisible; that mercantilist logic, when it enters intimate relational fields through bride-price, gold, and their contemporary equivalents, operates through three specific mechanisms---settlement, commodification, and indebtedness---that destroy the conditions for GRW generation; that the feminist political economy of social reproduction identifies the systematic non-recognition of care labour as a form of GRW extraction; and that the freedom-wealth dialectic is resolved not by choosing between material wealth and GRW but by organising material wealth as a condition for GRW generation.

The political economy framework cannot say: what GRW feels like from the inside; how the symbolic operations of the happiness jar make contact with the real register of the relational field; what the precise dynamical laws governing $\Delta$'s evolution are; or why the relational field generates happiness as its emergent signal rather than some other emergent property. These questions require the other frameworks that the paper has deployed.

\subsection{What the Formal Framework Can and Cannot Say}
\label{subsec:conc_formal}

\noindent
The formal framework---relational STDP, the structural remainder $\mathcal{R}(O(x), x) \neq 0$, the phase space account of GRW accumulation, the Relational Isomorphism Principle---can say: that GRW is stored in the structural modifications of the relational attractor landscape; that re-traversal produces additional holonomy that augments rather than depletes GRW; that the non-zero structural remainder is the formal basis of GRW's inexhaustibility; and that the Relational Isomorphism Principle provides a conjectural account of how symbolic operations in the happiness jar can produce real register effects through structural resonance.

The formal framework cannot say: whether the Relational Isomorphism Principle is correct, since this depends on the metaphysical commitment that the real is relational, which is a philosophical choice rather than a formally derivable theorem; what the precise form of the operator $O$ is, since the von Neumann problem for relational systems means that $O$ is not cleanly separable from the data it operates on; or whether the formal conjectures advanced here (the phase-space reconstruction conjecture, the parametric resonance analogy) will survive more rigorous formalisation. These are genuinely open questions, stated as such throughout the paper.

\subsection{What the Psychoanalytic Framework Can and Cannot Say}
\label{subsec:conc_psycho}

\noindent
The psychoanalytic framework can say: that traditional material wealth functions as a phallic signifier that organises intimate relations around the logic of having and lacking; that GRW is non-phallic in its structure, constitutively immune to the phallic logic of symbolic power; that GRW generation involves a circular drive of co-evolutionary generativity that is distinct from both desire and the death drive; and that the three pathologies of relational wealth---phallic substitution, imaginary fixation, and real register avoidance---provide a psychoanalytic typology of the ways in which the conditions for GRW generation can be destroyed while maintaining the appearance of relational flourishing.

The psychoanalytic framework cannot say: whether the concept of GRW as the real register's generative surplus is consistent with the full elaboration of Lacanian theory, since this requires a development of Lacanian theory that lies beyond the scope of this paper; or what the clinical implications of the GRW account are for psychoanalytic practice with intimate couples. These are open questions for psychoanalytic theory and practice.

\subsection{What the Feminist Framework Can and Cannot Say}
\label{subsec:conc_fem}

\noindent
The feminist framework can say: that the dominant model of subjectivity has systematically rendered invisible the care labour that generates GRW; that the care ethics tradition provides independent philosophical grounding for the GRW account of relational subjectivity; that GRW is non-phallic wealth whose cultivation requires and enables a genuinely relational subjectivity that is neither the liberal individual subject nor the care-ethical subject alone; and that the feminist political economy of GRW---the analysis of care labour as GRW production and its systematic non-recognition as a form of GRW extraction---provides the political economic content of the feminist critique of intimate relational wealth organisation.

The feminist framework cannot say: whether relational subjectivity as the GRW framework conceives it is compatible with the full range of feminist theoretical positions; how relational subjectivity is to be articulated in relation to the intersectional dimensions of gender, race, class, and sexuality that Crenshaw's framework requires; or what the implications of the GRW account are for non-binary and queer intimate relations, which the paper's focus on the heterosexual couple has inadequately addressed. These are significant limitations that future work must address.

\subsection{What the Ethical Framework Can and Cannot Say}
\label{subsec:conc_eth}

\noindent
The ethical framework can say: that GRW cultivation requires specific individual virtues, relational obligations, social ethical requirements, and political ethical demands that are dialectically unified in a generative spiral; that generative justice requires every party to have the right to generate relational value and to benefit from the GRW that their generative activity produces; that epistemic justice in intimate relational life requires epistemic symmetry, epistemic humility, and the active invitation of the other party's interpretive contribution; and that relational epistemic coercion---the imposition of one party's interpretive framework on the shared field---is both an epistemic injustice and a direct threat to GRW generation.

The ethical framework cannot say: how the relational consequentialist criterion is to be applied in specific cases where cultural context and universal ethical requirements are in tension; what the obligations of parties to intimate relations are when one party's investment in GRW generation is so asymmetric as to constitute a fundamental injustice that resists the usual relational remedies; or how the political ethical demands of the GRW framework are to be prioritised in conditions of political conflict and material scarcity. These are genuine ethical difficulties that admit of no general theoretical resolution.

\subsection{Three Concentric Circles}
\label{subsec:circles}

\noindent
The paper closes its conclusions with the same structural image that \pinkref{Paper~XIV} employed \citep{paperxiv}: three concentric circles that represent the three levels at which the GRW account operates.

The \emb{innermost circle} is the intimate relational field itself: the shared attractor landscape of a life built together, the accumulated holonomy of co-evolutionary engagement, the GRW stored in the geometry of the \emph{between}. This is the level of the happiness jar and the shared flower, of the wordless co-presence and the timely sharing, of the artistic processing and the re-traversal. It is the level that the phenomenological account of \pinkref{Paper~XIV} described and that the praxis account of the present paper has tried to make practically accessible.

The \emb{middle circle} is the political economy of intimate relational life: the distribution of care labour, the organisation of material wealth, the regulatory framework that either supports or undermines the conditions for GRW generation. This is the level of the feminist critique, the Hegelian-Marxist dialectic, the epistemic justice analysis. It is the level at which the individual practices of GRW cultivation are either supported or undermined by the social and economic structures within which intimate relational life unfolds.

The \emb{outermost circle} is the ontological and metaphysical framework that grounds the entire account: the relational ontology that holds that the real is relational in its basic structure, that the relational field is the primary unit of analysis, and that GRW is a genuine form of wealth that cannot be reduced to any of the categories that the dominant framework provides. This is the level at which the GRB series has been working throughout---the level of \pinkref{Paper~IX}'s holonomy framework \citep{paperix}, \pinkref{Paper~XIII}'s trust production mechanisms \citep{paperxiii}, and \pinkref{Paper~XIV}'s eudaimonistic reconstruction \citep{paperxiv}.

The three circles are not independent but mutually constitutive: the intimate relational field is embedded in the political economy that either supports or undermines it, and both are grounded in the ontological framework that makes sense of what they are and what they are for. A complete account of GRW must attend to all three circles simultaneously---which is why the paper has moved between the intimate and the structural, between the philosophical and the political economic, between the formal and the phenomenological, without settling permanently in any single register.

\subsection{Why This Paper Has No Synthetic Conclusion}
\label{subsec:nosynthesis}

\noindent
The reader who has reached this point may notice, as they noticed in \pinkref{Paper~XIV} \citep{paperxiv}, that the paper has not provided a synthetic conclusion: a paragraph that draws all the threads together into a single unified statement of what the paper has established. This absence is deliberate, and its reasons are the same as before.

A synthetic conclusion would imply that the paper has said the last word on its subject---that the concept of GRW is now complete, that the political economy of intimate relational wealth is now adequately theorised, that the practical account of how GRW is to be cultivated is now sufficient. But the paper has been honest throughout about what it does not know: the dynamical laws of $\Delta$, the origin of the operator $O$, the full implications of the Relational Isomorphism Principle, the application of the framework to non-heterosexual and intersectional intimate relations, the clinical implications for psychoanalytic practice, the specific policy interventions that the political ethical demands of the framework require. These are not loose ends to be tied up in a later paper; they are genuine open questions that the paper has generated, and they deserve to remain open.

The absence of a synthetic conclusion is also the paper's enactment of its own central thesis. GRW is not a stock that can be fully inventoried; it is a process that continues as long as co-evolutionary engagement continues. A paper about GRW that concluded with a comprehensive inventory of what has been established would be performing a settlement of what must remain generative. The paper that refuses synthetic closure is the paper that practises, in its own form, the non-possessive participation in the generative process that it recommends as the primary virtue of intimate relational life.

\zh{取之不尽，用之不竭。} Inexhaustible in the taking, inexhaustible in the using---because what is being taken and used is not a stock but a process, and the taking is itself a further moment of the process.
